% !TeX root = ../tfg.tex
% !TeX encoding = utf8


\iffalse
    ASUNCIONES CONDITIONAL MEAN INDEPENDENCE
    1. Introducción a la inferencia causal,
    1.1 las asunciones de o SUTVA, conditional mean independence,
    positivity y exchangeability y las potential outcomes
    y el estimado de la diferencia de riesgos y del odds ratio (OR).
    positivity
    exchangeability o SUTVA
    potential outcomes
    estimado de la diferencia de riesgos y del odds ratio (OR)
\fi


% TODO: find a better title?
\chapter{Framework de Potential Outcomes}

\section{Introducción}
Basándonos en \cite[capítulo 2]{targetedLearning}, consideremos el modelo estadístico $\mathcal{M} \ni P$,
donde $P$ es la distribución conjunta del triplete $(X, A, Y)$
de variables aleatorias observables, y supongamos que $Y$ es una variable aleatoria binaria.

A $X$ lo llamamos las covariantes, a $A$ el tratamiento, y a $Y$ el resultado o respuesta (\textit{outcome}, del inglés).

Consideramos entonces un modelo causal (\autoref{def:modelo}) $\mathcal{M}^*$
donde vive la distribución conjunta de un vector $(X, A, Y(0), Y(1))$,
$P^* \in \mathcal{M}^*$ donde $Y(0)$ e $Y(1)$ son dos variables aleatorias nuevas,
que en este momento no relacionamos con la variable $Y$ del modelo estadístico.
Como modelo causal, no es observable.

Veremos que $P^*$ determina $P$ en su totalidad bajo ciertos supuestos.

\begin{notac}
    Denotaremos $Y(0)$ e $Y(1)$ como las variables aleatorias correspondientes a un resultado
    sujeto al supuesto de que el tratamiento o intervención $A$ tenga valor $0$ o $1$,
    respectivamente, para \emph{toda} la población.

    Esto no se debe confundir con la variable estadística $Y$ condicionada al valor de $A$,
    puesto que el primero son variables en el modelo causal y el segundo lo son en el estadístico.

    Representan los resultados potenciales (\textit{Potential Outcomes}) que nos permitirán
    razonar de forma contrafactual.
\end{notac}

\section{Estimandos de interés} \label{sec:estimandos_de_interes}
En las siguientes definiciones consideramos una distribución $P^*$ de $(Y(0), Y(1), A)$
donde tanto $A$ como $Y(a)$ son variables binarias salvo mención de lo contrario.
$A$ es un factor de riesgo
o tratamiento
e $Y$ es la variable de respuesta.

% TODO: another(?) source
\begin{definicion}[Diferencia de Riesgo] \label{def:riskdiff} \cite[capítulo 1.3]{whatif}
    Definimos la diferencia de riesgo (RD por sus siglas en inglés) como:
    $$
        RD(P^*) = P^*(Y(1) = 1) - P^*(Y(0) = 1)
    $$
\end{definicion}

\begin{definicion}[ATE causal]\cite[sección 2.1]{mci}\label{def:causalATE}
    El ATE causal, del inglés \textit{Average Treatment Effect} o Efecto Medio de Tratamiento,
    es definido como
    $$ \psi^*(P^*) = \E_{P^*}(Y(1) - Y(0)) $$
    En este caso, la variable respuesta puede tomar varios valores discretos o ser continua.
    En el caso que $Y$ sea binaria, coincide con la diferencia de riesgo.
\end{definicion}

% TODO: check "covariantes" is a proper translation, and outcome thing
\iffalse
    \begin{observacion}[Interpretación] \cite[sección 1.3]{mci}
        El ATE mide la diferencia media entre los pacientes tratados y no tratados, ponderando
        por la predisposición a las covariantes en la población, $X$.
        % También depende de la distribución conjunta de $(Y, X, A)$.
        % En particular, el ATE depende de la distribución de covariantes $X$.
    \end{observacion}
\fi

\begin{definicion}[Riesgo Relativo] \label{def:relrisk} % TODO: source
    Definimos el Riesgo Relativo (RR) como:
    $$
        RR(P^*) = \frac{P^*(Y(1) = 1)}{P^*(Y(0) = 1)}
    $$
\end{definicion}

\begin{definicion}[Odds Ratio] \label{def:oddsratio} % TODO: source
    Definimos el Odds Ratio (OR) como:
    % OR = RR * \dfrac{P(\neg E|C)}{P(\neg E | I)} = RR * \dfrac{1 - P(E|C)}{1 - P(E|I)}
    \begin{align*}
        OR(P^*) & = \frac{P^*(Y(1) = 1)/P^*(Y(1) = 0)}{P^*(Y(0) = 1)/P^*(Y(0) = 0)}
        % OR & = \frac{P(Y(1) = 1)}{P(Y(0) = 1)} \cdot \frac{P(Y(0) = 0)}{P(Y(1) = 0)}
    \end{align*}
\end{definicion}

\begin{corolario}[Relación entre OR y RR]
    $$
        OR(P^*) = RR(P^*) \dfrac{P^*(Y(0) = 0)}{P^*(Y(1) = 0)}
    $$
\end{corolario}

\begin{corolario}[Caso de enfermedad rara] % TODO: prove
    % Si $P(Y = 0 | A = a) \approx 1 \; \forall a$, entonces $OR \approx RR$.
    Si $P^*(Y(a) = 0) \approx 1 \; \forall a$, entonces $OR \approx RR$.
\end{corolario}
\begin{proof}
    Tenemos que
    \[
        OR - RR = RR \cdot \left(\dfrac{P^*(Y(0) = 0)}{P^*(Y(1) = 0)} - 1\right)
    \]
    Bajo el límite $P^*(Y(a) = 0) \to 1$
    el último factor se anula.
    % Al realizar el experimento y aplicar consistencia \autoref{consistency},
    % la condición adicional $P(A = a) \nrightarrow 0$ será suficiente.
\end{proof}

Estudiaremos los estimandos estadísticos asociados a estos estimandos causales
en la \autoref{sec:identificacion_estimandos}.

\begin{ejemplo} % TODO: improve wording
    Un estudio médico puede estudiar la efectividad de un tratamiento en un
    grupo de pacientes de los cuales se registra la mortalidad en un plazo determinado.
    La variable binaria $A$ representa si se ha aplicado cierto tratamiento o no,
    la mortalidad es medida en la variable binaria $Y$,
    y además tenemos en cuenta una condición $X$ que puede afectar a las otras
    dos variables, como la presión arterial del paciente.
\end{ejemplo}


% \begin{definicion}[Independencia condicional]
% TODO: maybe put it in conceptos preliminares
% Sean $A$, $B$ y $C$ variables aleatorias.

% Se dice que $A$ y $B$ son condicionalmente independientes dado $C$, y se denota
% $ (A \perp B )| C $, si se verifica que
% $ P(A | B, C ) = P(A | C)$
% \end{definicion}



\iffalse
    \begin{ejemplo}
        Diferencia de riesgo: $Y_i(1) - Y_i(0)$,
        Riesgo relativo (RR): $Y_i(1) / Y_i(0)$,
        para un $i$ dado.
        % \item % TODO: does odd ratios fit here?
    \end{ejemplo}
\fi

\section{Supuestos del Potential Outcomes Framework}
\cite[capítulo 2.2]{mci}

\subsection{Motivación}
Dado un estimando causal $\psi^* : \mathcal{M}^* \to \R$, queremos encontrar
un estimando sobre la distribución observable que nos permita inferir sobre él,
usando que la distribución estadística se relaciona con la causal a través de
la observación: $P = \mathcal{O}(P^*)$.
Como los nombres indican, en particular querremos identificar el ATE causal con el ATE estadístico.

\begin{definicion}[Identificación]
    Diremos que un estimando causal $\varphi$ está identificado con un estimando estadístico $\psi$ cuando
    $ \varphi = \psi \circ \mathcal{O}$.
    En tal caso, denotaremos $\psi^* \equiv \varphi$.
\end{definicion}

Visualmente, queremos que el diagrama de la derecha sea conmutativo.
\nopagebreak

\begin{minipage}{0.45\textwidth}
    \centering
    \begin{align*}
        \mathcal{O}: \mathcal{M}^* \to & \; \mathcal{M} \\
        P^* \mapsto                    & \; P
    \end{align*}
\end{minipage}
\hfill
\begin{minipage}{0.45\textwidth}
    \centering
    $$
        \begin{tikzcd}
            \mathcal{M}^* \arrow{r}{{\mathcal{O}}} \arrow[swap]{dr}{\psi^*} & \mathcal{M} \arrow{d}{\psi} \\
            & \R^d
        \end{tikzcd}
    $$
\end{minipage}

% ( https://onlinelibrary.wiley.com/action/downloadSupplement?doi=10.1002%2Fsim.9234&file=sim9234-sup-0001-supinfo.pdf)
%
% https://www.jstor.org/stable/pdf/27590541.pdf?refreqid=fastly-default%3A596afd455befa7c2d5bf11298b9f67d5&ab_segments=&initiator=&acceptTC=1
% Los 3 primeros items son SUTVA TODO: write this out

% TODO: some intro


\subsection{Supuestos de identificación}

El primer problema que aparece es la mayor libertad que posee el modelo causal, que resulta en
que varias distribuciones causales tienen la misma distribución correspondiente en $\mathcal{M}$.
Eso es decir que $\mathcal{O}$ no es inyectiva, con lo que inmediatamente
las aplicaciones que llevan estimandos causales en estimandos estadísticos,
$\pi: \left\{\mathcal{M}^* \to \R \right\} \mapsto \left\{\mathcal{M} \to \R \right\}$ tampoco son inyectivas.

% TODO: thesaurus problema
Con el objectivo de lidiar con este problema se introducen las \textit{asunciones de identificación}.
Antes de ello, un ejemplo.

\begin{ejemplo}\cite[sección 2.2]{mci}
    Consideramos las tripletas $(A, Y(0),Y(1))$ con distribuciones en $\mathcal{M}^*$
    degeneradas tomando como valores posibles $(1, 0, 1)$ y $(1, -1, 1)$. Aplicando
    lo que llamaremos consistencia (\autoref{def:consistency}),
    tenemos que la distribución observable en ambos casos es la distribución degenerada de valor
    $(A, Y) = (1, 1)$.
    Este ejemplo específicamente pone de manifiesto la necesidad de la asunción de positividad
    (\autoref{def:positivity}),
    porque $A$ siempre toma el valor $1$ y por tanto siempre
    se da $Y = Y(1)$ mientras que el valor de $Y(0)$ resulta libre.
\end{ejemplo}

Los supuestos de identificación que hacemos son los siguientes, extraídos
de \cite[material suplementario S1]{luque2021} y \cite[capítulo 2.2]{mci}:

\begin{definicion}[Consistencia]
    \label{def:consistency}
    Una condición que vamos a imponer es la de \textbf{consistencia}, que enlaza
    las distribuciones causal y estadística:
    $$ Y = A Y(1) + (1 - A)Y(0) $$
    No confundir con el concepto de estimador estadístico consistente o asintóticamente insesgado,
    como definido en \autoref{def:insesgadez-asintotica}.
\end{definicion}

\begin{definicion}[Independencia condicional] \cite[1.1.5]{pearl2009causality}
    Se dice que dos variables $X, Y$ son independientes condicionadas a $Z$, denotado
    $X \indep Y | Z$, si
    \[
    P(X = x | Y = y, Z = z) = P(X = x | Z = z) \quad \forall x, y, z: P(Y = y, Z = z) > 0
    \]
    Esto es, conocida $Z$, conocer $Y$ no aporta información adicional sobre $X$.
\end{definicion}

\begin{definicion}[Intercambiabilidad condicional] \label{def:intercambiabilidad}
    También denominada independencia condicional, ignorabilidad o suficiencia causal.
    $$ Y(a) \indep A | X $$
    Esto es independencia entre $Y(a)$ y $A$, bajo condicionamiento a $X$. Debe leerse como
    $(Y(a) \indep A) | X$.
    % No confundir con indepedencia entre $Y(a)$ y $A | X$, denotada $ Y(a) \indep (A | X)$
    Una forma de verlo es que dado un paciente con ciertas covariantes $X$, la asignación al
    tratamiento no aporta información adicional.
    También significa que no hay más variables confusoras además de las representadas por $X$. % TODO: confusoras o simplemente covariantes?
    En particular, se obtiene
    $$
        \E[Y(a) | A = 1, X = x] = \E[Y(a) | A = 0, X = x] = \E[Y(a) | X = x] \;\forall a \in \{0, 1\}
    $$
\end{definicion}

\begin{definicion}[Positividad] \label{def:positivity}
    $$ P(A = a | X = x) > 0 \; \forall a \in \{0, 1\}, \; \forall x:  P(X = x) > 0$$
\end{definicion}

Dependiendo del autor la condición de consistencia vista en \ref{def:consistency}
se toma como una de ellas o no, puesto que su rol es de relacionar el modelo causal
y el modelo estadístico, y no constituye una restricción en la familia de modelos causales.

\begin{proposicion}
    Si el modelo causal y el estadístico son consistentes
    y $(X, A)$ sigue la misma distribución entre el modelo causal y el modelo estadístico:
    $$
        \begin{cases}
            P^*(X = x) = P(X = x) \;                 & \forall x    \\
            P^*(A = a | X = x) = P(A = a | X = x) \; & \forall x, a
        \end{cases}
    $$

    Entonces la distribución causal determina totalmente la distribución estadística por:
    $$ P(Y, A, X) = P^*([AY(1) + (1 - A)Y(0)] \; |\, A, X) P^*(A | X) P^*(X) $$
\end{proposicion}

\begin{proof} \cite[sección 2.1]{mci}
    % TODO:
    % Usando el teorema de la probablidad total:
    % \begin{align*}
    %     P(Y, A, X) &= P(X)P(A | X)P(Y | X, A) \\
    %                &= P^*(X)P^*(A | X)P(Y | X, A)
    % \end{align*}
    % y
    % \begin{equation*} %TODO: verify
    %     P^*(Y(0), Y(1), A, X) = P^*(X)P^*(A | X)P^*(Y(0), Y(1) | X, A)
    % \end{equation*}
    %
    % Usando la consistencia, debe verificarse que
    % \begin{equation*}
    %     P(Y | A, X) = P^*(AY(1) + (1-A) Y(0), A, X)
    % \end{equation*}
\end{proof}

La intercambiabilidad condicional está relacionada con el siguiente concepto:

\begin{definicion}[Intercambiabilidad marginal] \label{def:intercambiabilidad-marginal}
    Decimos que tenemos intercambiabilidad marginal o simplemente intercambiabilidad cuando
    se verifica
    $$ Y(a) \indep A \quad \forall a \in \{0, 1\} $$ % TODO: verify usage of a
    En el caso discreto:
    $$
        P(Y(a) = y | A = 1) = P(Y(a) = y | A = 0) = P(Y(a) = y) \quad \forall a \in \{0, 1\}, \; \forall y
    $$
    Eso es, el grupo tratado y el no tratado tienen la misma distribución de resultados potenciales.
    Por ello la nomenclatura: son intercambiables.
\end{definicion}

\begin{observacion}
    En el caso general de unas covariantes $X$ arbitrarias sin un grafo causal especificado
    como \autoref{subfig:observacional}, ninguna de las dos intercambiabilidades implica la otra.

    En un estudio observacional, con $X$ una causa común de $A$ e $Y$, puede darse
    la intercambiabilidad condicional y no darse la marginal.

    Si en cambio $X$ es un colisionador entre $A$ e $Y$,
    puede darse la intercambiabilidad marginal y no la condicional.

    Estas relaciones se estudian, entre otras formas, mediante criterios sobre grafos causales,
    como la $d$-separación \cite[1.2.3]{pearl2009causality} y el criterio de la puerta trasera
    o \emph{backdoor criterion} \cite[3.3.1]{pearl2009causality}.
\end{observacion}

\subsection{No interferencia y SUTVA}
\begin{definicion}[Unidad]\label{def:unidad}
    Llamaremos unidad a un único sujeto o un conjunto de sujetos que trataremos de forma colectiva.
\end{definicion}

\begin{ejemplo}
    En \cite{Haushofer2016}, trabajo asociado a GiveDirectly,
    se toma como unidad la unidad familiar para la mayoría de medidas socioeconómicas.
    Las medidas psicológicas sí consideran la unidad individual.
    Vale la pena observar que esto es a pesar de haber literatura mostrando que no lo son del todo,
    así reconocido en el propio artículo.
    Citando \cite{Haushofer2016}:
    \begin{quotation}
        % DONETODO: check this interpretation is actually what they meant (wording)
        In [the cited] literature, a common test of whether households are unitary is precisely the
        test of income pooling, which asks whether expenditure shares are different when money is received
        by husbands versus wives. Several papers have found this to be the case in observational data ([...]);
        in particular, female income is associated with larger expenditures on food and children.
    \end{quotation}
    Donde citan entre otros artículos, \cite{Thomas1990}, cuyo abstract empieza:
    \begin{quotation}
        If household income is pooled and then allocated to maximize welfare then income under the
        control of mothers and fathers should have the same impact on demand. With survey data on
        family health and nutrition in Brazil, the equality of parental income effects is rejected.
    \end{quotation}
\end{ejemplo}

\begin{definicion}[No interferencia]\label{def:noInterference}
    Se dice que se cumple la asunción de no interferencia cuando los efectos a nivel de unidad (\autoref{def:unidad}) son independientes.
    Es decir, que la exposición o no exposición de un sujeto al tratamiento no afecta al resultado de otros sujetos.
    % TODO: important: "no unobserved confounding" interpretation of this:
\end{definicion}


\begin{definicion}[SUTVA]
    \textit{Stable Unit Treatment Values Assumption} (\textit{SUTVA}) combina las condiciones de
    consistencia (\ref{def:consistency}) y de no interferencia (\ref{def:noInterference}).
    Ambas condiciones se satisfacen simultáneamente bajo muestreo IID.
\end{definicion}

\begin{ejemplo}[Violación de SUTVA]
    El efecto por el cual el tratamiento a unidad afecta a otras se le llama \textit{\say{spillover}} en inglés, o \say{derrame} en castellano.
    Un posible ejemplo de experimento donde no se cumple SUTVA es el análisis % por RCT % TODO: really por RCT?
    de los efectos de una transferencia incondicional de efectivo a algunas familias en una localidad africana.
    Aquí, la transferencia de efectivo no sólo afecta a la familia que lo recibe, sino
    que también puede afectar a la localidad por el incremento en el consumo
    de esta familia.
    Este es un ejemplo del trabajo que hacen en organizaciones como GiveDirectly.
    En \cite{Haushofer2016}, seguido posteriormente por \cite{Egger2022},
    dividen el grupo de control entre \say{Control puro}, para familias no afectadas por
    el efecto de spillover, y el grupo de \say{Control con spillover}, que incluye familias
    de la misma localidad que los recipientes.
    Se puede encontrar una explicación más divulgativa de estos artículos en \cite{GiveDirectly2016}.

    Otro ejemplo es un estudio de la efectividad de una vacuna donde el muestreo se asigna por grupos (como un grupo
    de estudiantes de la misma clase)\cite{oxChapterCausal}.
\end{ejemplo}
